% Appendix (supplementary material).  v2, 13 Oct 2026 (FINAL_SUBMISSION_REVIEW_ASTRA: implementation details, full baseline
% tables, per-seed tables, protocols and extensions moved here from the main text; Session-9 scope corrected).
% Tables generated by manuscript/make_tables.py (except tab_related, hand-written).
\clearpage
\appendix
\section{Contained classical cycles, implementation details and numerical checks}\label{app:impl}
\begin{table}[tbp]\centering\footnotesize\setlength{\tabcolsep}{4pt}
\caption{Classical cycles contained in the architecture (Proposition~\ref{prop:containment}, same hierarchy, coarse solution, deflation and Galerkin output; $\nu=4$): fixed-interval Chebyshev smoothing ($\lambda=3$, same weights on every level) and weighted Jacobi ($\omega=2/3$), together with the adaptive Chebyshev reference used in the other tables. The contained cycles are evaluated in double precision with the unregularised Galerkin solves of the mathematical construction; the regularisation and finite-precision safeguards of the deployed network are described in Section~\ref{sec:properties}. (a) Mean energy errors and paired $G$ of each contained cycle over the reference. (b) Paired $G$ of each classical cycle over OA-VCB ([95\% CI] for the contained cycles) and OA-VCB's win fraction $W$. Elasticity: no deflation; one data set.}\label{tab:contained}
{\small (a) Mean energy errors}\par\smallskip
\begin{adjustbox}{max width=\linewidth}
\begin{tabular}{lllllll}\toprule
set & OA-VCB & reference & fixed Cheb. & Jacobi & fixed Cheb./ref.\ $G$ & Jacobi/ref.\ $G$ \\\midrule
channels, seed 1 & 2.55e-3 & 2.59e-2 & 2.58e-2 & 4.03e-2 & 0.995 & 1.57 \\
channels, seed 2 & 2.79e-3 & 2.50e-2 & 2.49e-2 & 3.85e-2 & 0.995 & 1.55 \\
channels, seed 3 & 2.92e-3 & 2.63e-2 & 2.62e-2 & 4.05e-2 & 0.995 & 1.54 \\
barriers, seed 1 & 4.36e-2 & 7.51e-2 & 7.62e-2 & 8.63e-2 & 1.052 & 1.51 \\
barriers, seed 2 & 6.96e-2 & 8.51e-2 & 8.61e-2 & 9.64e-2 & 1.040 & 1.45 \\
barriers, seed 3 & 6.07e-2 & 8.33e-2 & 8.50e-2 & 9.50e-2 & 1.056 & 1.50 \\
$k_0=0$, seed 1 & 7.62e-3 & 8.52e-2 & 8.84e-2 & 1.15e-1 & 1.039 & 1.35 \\
$k_0=0$, seed 2 & 7.58e-3 & 8.04e-2 & 8.36e-2 & 1.09e-1 & 1.041 & 1.36 \\
$k_0=0$, seed 3 & 8.13e-3 & 8.55e-2 & 8.85e-2 & 1.15e-1 & 1.036 & 1.35 \\
L4, $k_0=3$, seed 1 & 2.21e-3 & 2.01e-2 & 2.02e-2 & 3.08e-2 & 1.008 & 1.54 \\
elasticity (prelim.) & 2.47e-2 & 5.50e-2 & 5.38e-2 & 9.84e-2 & 0.975 & 1.82 \\
\bottomrule\end{tabular}
\end{adjustbox}
\par\medskip
{\small (b) Gains of OA-VCB}\par\smallskip
\begin{adjustbox}{max width=\linewidth}
\begin{tabular}{llllll}\toprule
set & ref./OA-VCB $G$ & fixed Cheb./OA-VCB $G$ & $W$ & Jacobi/OA-VCB $G$ & $W$ \\\midrule
channels, seed 1 & 10.23 & 10.19 [9.69, 10.66] & 1.00 & 16.02 [15.27, 16.73] & 1.00 \\
channels, seed 2 & 9.76 & 9.71 [9.12, 10.27] & 1.00 & 15.11 [14.16, 16.00] & 1.00 \\
channels, seed 3 & 9.68 & 9.63 [9.07, 10.19] & 1.00 & 14.95 [14.08, 15.79] & 1.00 \\
barriers, seed 1 & 2.13 & 2.24 [1.99, 2.51] & 0.93 & 3.22 [2.83, 3.64] & 0.94 \\
barriers, seed 2 & 1.93 & 2.01 [1.76, 2.28] & 0.83 & 2.80 [2.39, 3.27] & 0.87 \\
barriers, seed 3 & 1.93 & 2.04 [1.78, 2.32] & 0.90 & 2.90 [2.53, 3.33] & 0.93 \\
$k_0=0$, seed 1 & 11.39 & 11.83 [11.38, 12.30] & 1.00 & 15.34 [14.74, 15.93] & 1.00 \\
$k_0=0$, seed 2 & 10.91 & 11.36 [10.78, 11.91] & 1.00 & 14.80 [14.05, 15.52] & 1.00 \\
$k_0=0$, seed 3 & 10.82 & 11.21 [10.77, 11.67] & 1.00 & 14.57 [14.01, 15.17] & 1.00 \\
L4, $k_0=3$, seed 1 & 9.72 & 9.80 [9.13, 10.42] & 1.00 & 14.98 [13.93, 15.92] & 1.00 \\
elasticity (prelim.) & 2.32 & 2.26 [2.15, 2.37] & 1.00 & 4.22 [3.93, 4.53] & 1.00 \\
\bottomrule\end{tabular}
\end{adjustbox}
\end{table}

\paragraph{Contained classical cycles.} Table~\ref{tab:contained} evaluates the two contained members of
Proposition~\ref{prop:containment} on the main test sets. The contained fixed-interval Chebyshev member is as accurate as the
adaptive reference to within a few per cent (paired ratio $0.975$--$1.056$), so OA-VCB's margins over it are essentially those
over the reference ($G = 9.6$--$10.2$ on channels, $2.0$--$2.2$ on barriers, $2.26$ for elasticity). The contained Jacobi member
is less accurate on every set ($96$--$100\%$ of the samples on the scalar sets).

\paragraph{Numerical check of the containment.} With the weights of the construction in Proposition~\ref{prop:containment},
the network and the classical cycle agree to $3.6\times10^{-10}$ for Jacobi weights (diffusion, \path{mathcheck/check_oa_vcp.log};
$1.1\times10^{-10}$ tensor coefficient, \path{mathcheck/check_oa_vcb_tensor.log}; $1.4\times10^{-10}$ elasticity,
\path{mathcheck/check_oa_vcb_elasticity.log}; $\omega = 2/3$, $\nu = 4$), to $4.8\times10^{-10}$ for the Chebyshev weights of (b)
with $\lambda = 3$ and $q = 10$, and to $3.7\times10^{-10}$ for the first construction with a random cubic $p$ whose residual
polynomial has a pair of complex roots (\path{mathcheck/check_poly_containment.log}; maximum relative differences of the
outputs on twelve barrier test samples). On every level operator of these samples the product form of (b) and the three-term
recurrence of the Chebyshev semi-iteration give the same smoother vectors to a relative difference of $5.9\times10^{-16}$.

\paragraph{Transposed transfers.} With $T = [\Phi_k^\star\;E_F]$, \eqref{eq:inverse} reads
$A_k^{-1} = T\,\operatorname{diag}(S_k,A_{FF})^{-1}\,T^\top$: the residual is analysed in the \emph{dual} coordinates $T^\top r$
and the correction synthesised with $T$. Any invertible synthesis $T$ admits this form. What it does not admit is a
\emph{block-diagonal} latent map between the primal pair $T^{-1}$ (analysis) and $T$ (synthesis), as used by exactly
invertible lifting networks: there the latent map would have to equal $T^{-1}A_k^{-1}T = \operatorname{diag}(S_k,A_{FF})^{-1}\,T^\top T$,
which is block diagonal only if $T^\top T$ is (\path{mathcheck/check_oa_transpose_pair.log}). With $X^{(m)}$ replaced by the
exact $-A_{FF}^{-1}A_{FC}$ and exact level solves the cycle reproduces $A^{-1}$ to $10^{-12}$--$10^{-14}$; with the approximation
the coupling between levels that $X^{(m)}$ leaves is handled by the multiplicative structure of the cycle.

\paragraph{Reproduction of the kernel (Proposition~\ref{prop:constants}).} On level $k$ the statement covers the fine nodes
whose $m$-step neighbourhood in the level-$k$ (Galerkin) graph satisfies both hypotheses. On the finest level that
neighbourhood spans $m$ mesh widths, and the measured rigid-body defect of the elasticity transfers is at round-off for most
rows (median $7\times10^{-17}$), but not for all (maximum $0.40$ over rows not directly coupled to $\Gamma_D$, whose $m$-step
neighbourhoods may still reach it). On coarser levels the Galerkin stencils widen, the hypotheses cover fewer rows (41\% on the
second level, none on the third) and the defects grow (median over all rows $1.5\times10^{-4}$ and $1.6\times10^{-2}$;
\path{mathcheck/check_oa_vcb_elasticity.log}). We report these measured defects rather than a physical-width bound.

\paragraph{Homogeneity in finite precision.} In the single-precision implementation $\sigma$ is evaluated as a scaled root
mean square ($\max|z|\cdot\|z/\max|z|\|_{\mathrm{rms}}$, which cannot underflow), and $\mathcal{N}(\lambda r) = \lambda\mathcal{N}(r)$
holds to $7\times10^{-7}$ for residual scales from $10^{-20}$ to $10^{8}$ on CPU and GPU (\path{mathcheck/check_projection_precision.log}).
At $10^{-30}$ the smallest output components fall below the smallest normal single-precision number; GPU kernels flush such
values to zero and the deviation grows to $3\times10^{-5}$ (finiteness and the energy bound are unaffected).

\paragraph{The energy safeguard in the implementation.} The implementation differs from Proposition~\ref{prop:safeguard} in
three ways; each preserves energy non-increase, none preserves the minimisation equality in general. (i)~\emph{Data.}
Deflation and the output projection use the residual and the fine-level operator in double precision, as assembled; the
cycle itself runs in single precision, which changes the directions $W$ but not the argument. (ii)~\emph{Regularisation.} The
Jacobi-scaled Gram matrix $G$ is symmetrised and a ridge $R = 10^{-10}I$ is added; for $(G+R)y = g$ the energy change is
$-y^\top Gy - 2y^\top Ry\le 0$, but $y$ is no longer the minimiser over $\operatorname{span}W$ when $R$ matters (nearly
dependent directions). (iii)~\emph{Acceptance.} For the computed correction $d$ of a residual $b$ (the deflation step and the
output projection alike), energy non-increase is equivalent to $2d^\top b - d^\top A d\ge 0$; this is evaluated with the
double-precision operator, and if it fails $d$ is replaced by $d\,\max(d^\top b,0)/d^\top Ad$, which may shrink or reject it.
The implemented guarantee is therefore $\|e-c\|_A\le\|e - x_0\|_A\le\|e\|_A$ up to double-precision rounding in the acceptance
test, with $c$ not necessarily the minimiser. No eigendecomposition is used in the forward pass, so the gradients stay defined
at repeated eigenvalues and rank changes; zero right-hand sides are mapped to zero explicitly. Numerically
(\path{mathcheck/check_projection_precision.log}): no violation at tolerance $10^{-12}$ on real samples with random and
classical weights, on 210 stress cases with nearly or exactly duplicate directions (as projection and as deflation with zero
padding), or for residual scales from $10^{-30}$ to $10^{8}$; finite gradients at duplicate columns and at zero and tiny
residuals; the acceptance step never triggered. An earlier implementation formed $G$ in single precision and broke the bound in
11 of the 210 stress cases. Re-evaluating trained checkpoints with the present implementation changed the corrections by at most
$7\times10^{-6}$ relative and the mean errors by less than $10^{-5}$ relative, with no acceptance failures and no rank loss
(\path{diagnostics_cost/checkpoint_audit_*.log}).

\section{Baselines: sources, adaptations and full results}\label{app:baselines}
\paragraph{Sources and adaptations.} All baselines are trained with the protocol of Section~\ref{sec:protocol} (same data,
same base loss, optimiser, schedule and epochs; the projected variants add the plain-output term of Section~\ref{sec:output})
and start from $u_c$, which is added back to their output; they predict only the correction. Except for the two excluded
published-size runs below, the width of each model is the largest integer, found by bisection, for which the model has at most
$22{,}000$ parameters; no other hyperparameter was tuned per class.
\begin{itemize}
\item \emph{OA-Lift} (\texttt{oa-ll-op}): our earlier operator-adapted lifting neural operator on the same mesh hierarchy, with
the error-equation inputs and operator weights of OA-VCB.
\item \emph{GNN} and \emph{GNN-op}: an encode--process--decode message-passing network of the MeshGraphNet type
\cite{pfaff2021} on the edges of the finest mesh, ten message-passing steps, re-implemented; GNN-op adds the operator weights
$A_{ij}/A_{ii}$ as edge features and the scaled residual as node input.
\item \emph{Transolver} and \emph{Transolver-op}: re-implementation of the Physics-Attention architecture of \cite{wu2024transolver}
(soft assignment of mesh points to learnable slices, attention among slice tokens), acting directly on the mesh nodes.
\item \emph{FNO-op}: a simplified four-layer Fourier neural operator in the spirit of \cite{li2021fno} on the $63\times63$
grid (our implementation, $\min(12, n_g/2)$ modes per direction). It differs from the published FNO architecture in three
respects: each spectral convolution retains only the low-frequency corner $[0,m)\times[0,m)$ of the real FFT, so the modes
with negative first wavenumber, which the published layer also weights, are discarded; the grid is not padded before the FFT;
and a GELU follows every Fourier layer, including the last, before a two-layer pointwise output map. The FNO-op rows therefore
measure this variant, not the published FNO.
\item \emph{WNO-op}: a light wavelet neural operator layer in the spirit of \cite{tripura2023wno} (our implementation: four
layers, two-level db4 transform with periodised boundary handling, learnable mixing per subband, spatial skip). The channel rows
use the \texttt{pytorch\_wavelets} backend, the barrier rows a fallback implementation of the same periodised db4 transform.
\item \emph{LiNO-op}: re-implementation of the lifting-based multiresolution neural operator of \cite{pandey2026lino}
(Algorithm~1 and Section~2.3 of that paper: convolutional encoder, three learnable lifting levels, coarse and directional
residual operators shared across levels, exact inverse lifting, convolutional decoder); grids are zero-padded to a multiple of
$2^3$ and cropped back.
\item \emph{MgNO} and \emph{MgNO-op}: the model classes of the official MgNO code \cite{he2024mgno} (repository
\url{https://github.com/xlliu2017/MgNO}, MIT licence) unchanged except for a removed import, wrapped for our channel-last grid
interface; grids are padded to a multiple of $2^{\ell-1}$ for $\ell$ levels and cropped back; default V-cycle of the authors'
Darcy script.
\item Excluded runs (Table~\ref{tab:excluded}): \emph{WNO, published architecture}, a port of version 2.0.0 of the authors'
repository (db6 transform via \texttt{pytorch\_wavelets}, spatially varying weights, published width), and \emph{MgNO,
published size} (official code, width 24), both trained without the parameter cap.
\end{itemize}
All grid models receive their inputs interpolated piecewise linearly from the mesh to the grid, with a domain-mask channel,
and return their output bilinearly interpolated to the mesh nodes (\appref{app:bridge}). The ``-op'' variants receive the
scaled residual $D^{-1}r/\sigma$ in place of the $u_c$ input channel and return $\sigma$ times their output.

\paragraph{Full results.} Tables~\ref{tab:channels} and~\ref{tab:barriers} list all baselines and further classical rows,
Tables~\ref{tab:channelsfull} and~\ref{tab:barriersfull} their per-seed means, and Table~\ref{tab:perseed} the per-seed paired
values of OA-VCB and OA-VC. On channels, OA-Lift is $6\%$ more accurate than the reference on average and wins on $63\%$ of the
samples; GNN-op, GNN and Transolver-op are less accurate than the reference by factors of $1.2$, $1.7$ and $7.1$, and the grid
models by factors of $17$--$25$. Relative to GNN-op, the grid models' nodal errors ($4.7$--$8.0\times10^{-2}$) are $3$--$6$ times
larger and their energy errors $15$--$22$ times larger; this is consistent with the bridge handicap, but Transolver-op, which
uses no bridge, shows a weaker form of it (nodal $5.3\times10^{-2}$, energy $0.185$), so the bridge is not its only cause. On
barriers none of the baselines comes within a factor of $6$ of the reference, and four of the five grid models are on average at
or above the error of $u_c$.
\begin{table}[tbp]\centering\footnotesize\setlength{\tabcolsep}{4pt}
\caption{Channels (contrast $10^4$, L3, $k_0=2$), held-out data seeds 1--3, 100 test operators each. Energy error: mean and sample s.d.\ of the three per-seed means; $G$: paired geometric-mean ratio of the classical reference's error to the model's over the pooled 300 samples; $W$: fraction of samples on which the model is more accurate than the reference. $^\dagger$ Grid models, through the grid bridge (Section~\ref{sec:baselines}).}\label{tab:channels}
\begin{adjustbox}{max width=\linewidth}
\begin{tabular}{lrllll}\toprule
model & params & energy (mean $\pm$ s.d.) & nodal rel.\ $\ell^2$ & $G$ [95\% CI] & $W$ \\\midrule
OA-VCB (ours) & 20{,}672 & 2.75e-3 $\pm$ 1.85e-4 & 8.76e-4 & 9.89 [9.57, 10.20] & 1.00 \\
OA-VC & 20{,}672 & 3.46e-3 $\pm$ 5.44e-4 & 1.09e-3 & 8.80 [8.45, 9.12] & 1.00 \\
OA-Lift & 21{,}603 & 2.39e-2 $\pm$ 3.69e-4 & 1.01e-2 & 1.06 [1.04, 1.09] & 0.63 \\
GNN-op & 20{,}273 & 3.01e-2 $\pm$ 6.76e-4 & 1.46e-2 & 0.84 [0.82, 0.86] & 0.19 \\
GNN & 20{,}257 & 4.21e-2 $\pm$ 1.28e-3 & 1.60e-2 & 0.60 [0.58, 0.61] & 0.01 \\
Transolver-op & 19{,}521 & 1.85e-1 $\pm$ 2.57e-3 & 5.30e-2 & 0.14 [0.13, 0.14] & 0.00 \\
FNO-op$^\dagger$ & 18{,}561 & 6.57e-1 $\pm$ 5.42e-2 & 8.03e-2 & 0.04 [0.04, 0.04] & 0.00 \\
WNO-op$^\dagger$ & 21{,}377 & 4.56e-1 $\pm$ 8.01e-3 & 5.16e-2 & 0.06 [0.06, 0.06] & 0.00 \\
LiNO-op$^\dagger$ & 19{,}482 & 6.08e-1 $\pm$ 2.45e-2 & 4.67e-2 & 0.04 [0.04, 0.05] & 0.00 \\
MgNO$^\dagger$ & 21{,}140 & 5.76e-1 $\pm$ 4.04e-2 & 4.95e-2 & 0.05 [0.04, 0.05] & 0.00 \\
MgNO-op$^\dagger$ & 21{,}140 & 5.60e-1 $\pm$ 3.48e-3 & 4.84e-2 & 0.05 [0.05, 0.05] & 0.00 \\
\midrule
classical reference & -- & 2.57e-2 $\pm$ 6.95e-4 & 1.19e-2 & 1 & -- \\
classical V(4,4), no defl./proj. & -- & 5.36e-2 $\pm$ 2.04e-3 & 3.31e-2 & 0.48 [0.47, 0.48] & 0.00 \\
two-grid PCG, 4 its & -- & 1.14e-2 $\pm$ 3.47e-4 & 3.58e-3 & 2.34 [2.29, 2.40] & 1.00 \\
two-grid PCG, 6 its & -- & 2.08e-3 $\pm$ 2.06e-5 & 5.97e-4 & 14.13 [13.49, 14.81] & 1.00 \\
two-grid PCG, 8 its & -- & 3.78e-4 $\pm$ 2.45e-5 & 1.08e-4 & 89.30 [83.44, 95.72] & 1.00 \\
$u_c$ alone & -- & 3.73e-1 $\pm$ 5.82e-3 & 1.44e-1 & 0.07 [0.07, 0.07] & 0.00 \\
\bottomrule\end{tabular}
\end{adjustbox}
\end{table}

\begin{table}[tbp]\centering\footnotesize\setlength{\tabcolsep}{4pt}
\caption{Barriers (contrast $10^{-4}$, source only where $a\ge1/2$, L3, $k_0=2$), held-out data seeds 1--3. Columns as in Table~\ref{tab:channels}. The barrier WNO-op runs used the fallback discrete wavelet transform (\appref{app:baselines}).}\label{tab:barriers}
\begin{adjustbox}{max width=\linewidth}
\begin{tabular}{lrllll}\toprule
model & params & energy (mean $\pm$ s.d.) & nodal rel.\ $\ell^2$ & $G$ [95\% CI] & $W$ \\\midrule
OA-VCB (ours) & 20{,}672 & 5.80e-2 $\pm$ 1.32e-2 & 3.32e-2 & 2.00 [1.86, 2.14] & 0.88 \\
OA-VC & 20{,}672 & 3.73e-1 $\pm$ 1.56e-1 & 4.10e-1 & 0.61 [0.54, 0.70] & 0.43 \\
OA-Lift & 21{,}603 & 3.52e-1 $\pm$ 3.37e-2 & 3.78e-1 & 0.15 [0.14, 0.16] & 0.00 \\
GNN-op & 20{,}273 & 3.61e-1 $\pm$ 1.75e-2 & 4.08e-1 & 0.12 [0.11, 0.13] & 0.00 \\
GNN & 20{,}257 & 4.44e-1 $\pm$ 2.27e-2 & 4.09e-1 & 0.08 [0.07, 0.09] & 0.00 \\
Transolver-op & 19{,}521 & 3.79e-1 $\pm$ 2.48e-2 & 4.09e-1 & 0.11 [0.10, 0.12] & 0.00 \\
FNO-op$^\dagger$ & 18{,}561 & 4.76e-1 $\pm$ 2.42e-2 & 4.13e-1 & 0.08 [0.07, 0.09] & 0.00 \\
WNO-op$^\dagger$ & 21{,}377 & 5.13e-1 $\pm$ 4.82e-2 & 4.06e-1 & 0.08 [0.07, 0.09] & 0.00 \\
LiNO-op$^\dagger$ & 19{,}482 & 4.67e-1 $\pm$ 2.02e-2 & 4.13e-1 & 0.08 [0.07, 0.09] & 0.00 \\
MgNO$^\dagger$ & 21{,}140 & 4.28e-1 $\pm$ 2.42e-2 & 4.11e-1 & 0.08 [0.08, 0.09] & 0.00 \\
MgNO-op$^\dagger$ & 21{,}140 & 5.38e-1 $\pm$ 7.51e-2 & 4.05e-1 & 0.08 [0.07, 0.08] & 0.00 \\
\midrule
classical reference & -- & 8.11e-2 $\pm$ 5.34e-3 & 6.86e-2 & 1 & -- \\
classical V(4,4), no defl./proj. & -- & 2.67e-1 $\pm$ 3.37e-2 & 3.20e-1 & 0.42 [0.39, 0.44] & 0.00 \\
two-grid PCG, 4 its & -- & 1.34e-1 $\pm$ 2.10e-2 & 1.34e-1 & 1.67 [1.48, 1.89] & 0.72 \\
two-grid PCG, 6 its & -- & 5.05e-2 $\pm$ 2.24e-2 & 4.24e-2 & 10.36 [8.76, 12.23] & 0.89 \\
two-grid PCG, 8 its & -- & 2.43e-2 $\pm$ 1.66e-2 & 1.91e-2 & 68.20 [55.32, 83.79] & 0.96 \\
$u_c$ alone & -- & 4.59e-1 $\pm$ 2.23e-2 & 4.25e-1 & 0.07 [0.07, 0.08] & 0.00 \\
\bottomrule\end{tabular}
\end{adjustbox}
\end{table}

\begin{table}[tbp]\centering\footnotesize\setlength{\tabcolsep}{4pt}
\caption{Per-seed mean energy errors and the number of test samples on which the model is worse in energy than $u_c$ (seeds 1 / 2 / 3); channels (contrast $10^4$, l3, $k_0=2$), held-out data seeds 1--3, 100 test operators each.}\label{tab:channelsfull}
\begin{adjustbox}{max width=\linewidth}
\begin{tabular}{lllll}\toprule
model & seed 1 & seed 2 & seed 3 & worse than $u_c$ \\\midrule
OA-VCB (ours) & 2.55e-3 & 2.79e-3 & 2.92e-3 & 0 / 0 / 0 \\
OA-VC & 4.08e-3 & 3.10e-3 & 3.19e-3 & 0 / 0 / 0 \\
OA-Lift & 2.42e-2 & 2.35e-2 & 2.41e-2 & 0 / 0 / 0 \\
GNN-op & 2.99e-2 & 2.96e-2 & 3.09e-2 & 0 / 0 / 0 \\
GNN & 4.10e-2 & 4.18e-2 & 4.35e-2 & 0 / 0 / 0 \\
Transolver-op & 1.85e-1 & 1.82e-1 & 1.87e-1 & 0 / 0 / 0 \\
FNO-op$^\dagger$ & 6.25e-1 & 7.20e-1 & 6.27e-1 & 96 / 96 / 93 \\
WNO-op$^\dagger$ & 4.47e-1 & 4.62e-1 & 4.59e-1 & 65 / 68 / 72 \\
LiNO-op$^\dagger$ & 5.95e-1 & 5.93e-1 & 6.36e-1 & 96 / 91 / 92 \\
MgNO$^\dagger$ & 6.20e-1 & 5.66e-1 & 5.41e-1 & 89 / 92 / 80 \\
MgNO-op$^\dagger$ & 5.57e-1 & 5.58e-1 & 5.64e-1 & 87 / 91 / 91 \\
\bottomrule\end{tabular}
\end{adjustbox}
\end{table}

\begin{table}[tbp]\centering\footnotesize\setlength{\tabcolsep}{4pt}
\caption{Per-seed mean energy errors and the number of test samples on which the model is worse in energy than $u_c$ (seeds 1 / 2 / 3); barriers (contrast $10^{-4}$, source only where $a\ge1/2$, l3, $k_0=2$), held-out data seeds 1--3.}\label{tab:barriersfull}
\begin{adjustbox}{max width=\linewidth}
\begin{tabular}{lllll}\toprule
model & seed 1 & seed 2 & seed 3 & worse than $u_c$ \\\midrule
OA-VCB (ours) & 4.36e-2 & 6.96e-2 & 6.07e-2 & 0 / 0 / 0 \\
OA-VC & 2.33e-1 & 5.41e-1 & 3.45e-1 & 7 / 7 / 12 \\
OA-Lift & 3.19e-1 & 3.86e-1 & 3.51e-1 & 1 / 6 / 2 \\
GNN-op & 3.43e-1 & 3.78e-1 & 3.63e-1 & 0 / 0 / 0 \\
GNN & 4.20e-1 & 4.65e-1 & 4.48e-1 & 12 / 21 / 1 \\
Transolver-op & 3.53e-1 & 4.02e-1 & 3.82e-1 & 0 / 0 / 0 \\
FNO-op$^\dagger$ & 4.52e-1 & 5.00e-1 & 4.77e-1 & 18 / 24 / 17 \\
WNO-op$^\dagger$ & 4.78e-1 & 5.68e-1 & 4.93e-1 & 17 / 20 / 17 \\
LiNO-op$^\dagger$ & 4.56e-1 & 4.91e-1 & 4.56e-1 & 20 / 18 / 21 \\
MgNO$^\dagger$ & 4.02e-1 & 4.49e-1 & 4.33e-1 & 3 / 10 / 6 \\
MgNO-op$^\dagger$ & 4.92e-1 & 6.25e-1 & 4.97e-1 & 33 / 26 / 27 \\
\bottomrule\end{tabular}
\end{adjustbox}
\end{table}

\begin{table}[tbp]\centering\footnotesize\setlength{\tabcolsep}{4pt}
\caption{Per-seed results for OA-VCB, its classical reference and OA-VC (mean energy errors; reference/model $G$ [95\% CI] and $W$).}\label{tab:perseed}
\begin{adjustbox}{max width=\linewidth}
\begin{tabular}{llllllll}\toprule
class & seed & OA-VCB & reference & $G$ [95\% CI] & $W$ & OA-VC & $G$, $W$ \\\midrule
channels & 1 & 2.55e-3 & 2.59e-2 & 10.23 [9.75, 10.71] & 1.00 & 4.08e-3 & 8.83, 0.99 \\
channels & 2 & 2.79e-3 & 2.50e-2 & 9.76 [9.16, 10.31] & 1.00 & 3.10e-3 & 8.70, 1.00 \\
channels & 3 & 2.92e-3 & 2.63e-2 & 9.68 [9.12, 10.24] & 1.00 & 3.19e-3 & 8.88, 1.00 \\
\midrule
barriers & 1 & 4.36e-2 & 7.51e-2 & 2.13 [1.90, 2.38] & 0.92 & 2.33e-1 & 0.71, 0.44 \\
barriers & 2 & 6.96e-2 & 8.51e-2 & 1.93 [1.70, 2.18] & 0.83 & 5.41e-1 & 0.65, 0.48 \\
barriers & 3 & 6.07e-2 & 8.33e-2 & 1.93 [1.70, 2.19] & 0.89 & 3.45e-1 & 0.51, 0.36 \\
\bottomrule\end{tabular}
\end{adjustbox}
\end{table}

\begin{table}[tbp]\centering\footnotesize\setlength{\tabcolsep}{4pt}
\caption{Published-size grid models trained under the protocol (channels, seed 1, same loss and schedule, no parameter cap) and excluded from superiority claims: mean energy error, nodal relative $\ell^2$ error, and samples worse than $u_c$ in energy.}\label{tab:excluded}
\begin{adjustbox}{max width=\linewidth}
\begin{tabular}{llllll}\toprule
model & implementation & params & energy & nodal rel.\ $\ell^2$ & worse than $u_c$ \\\midrule
WNO{,} published architecture & Tripura--Chakraborty v2.0.0 port{,} pytorch\_wavelets & 74{,}865 & 1.54e0 & 1.43e-1 & 100/100 \\
MgNO{,} published size & official model code{,} width 24 & 658{,}020 & 3.76e-1 & 1.46e-1 & 0/100 \\
$u_c$ alone (reference) & -- & -- & 3.76e-1 & 1.46e-1 & -- \\
\bottomrule\end{tabular}
\end{adjustbox}
\end{table}

\paragraph{Excluded runs.} Table~\ref{tab:excluded} lists the two published-size grid models that did not train successfully
under the protocol (training loss $1.000$ for MgNO, $12.8$ for WNO at the last logged epochs); they are excluded from
superiority claims. The representation floor alone (\appref{app:bridge}) does not explain these failures: a field
representable through the bridge attains a training-loss value of about $0.25$, below that of the zero output, $1$. This does not
exclude optimisation difficulties induced by the bridge; the cause was not established. The failures say nothing about these
architectures on the uniform-grid problems they were designed for.

\section{Baselines on the published Darcy benchmark}\label{app:published}
% Session 9 (12 Oct 2026): kaggle_runs/published_darcy.py; records Results_kaggle_session 9/; table generated by make_tables.py.
To check several baseline implementations on the problem they were designed for, we train them on the piecewise-constant
Darcy benchmark of \cite{li2021fno}: the files \texttt{piececonst\_r421\_N1024\_smooth1/2}, subsampled by five to
$85\times85$, the first $1000$ samples of the first file for training and the first $100$ of the second for testing, $500$
epochs, mean relative $L^2$ error on the test set, and the optimiser settings of the respective papers, at their published
sizes (Table~\ref{tab:published}; script \path{kaggle_runs/published_darcy.py}). Each value comes from one training run. The
published values on this benchmark are $1.08\%$ for FNO and $0.84\%$ for WNO \cite{tripura2023wno}.
\begin{table}[tbp]\centering\footnotesize\setlength{\tabcolsep}{4pt}
\caption{Baselines on the published piecewise-constant Darcy benchmark of \cite{li2021fno} ($85\times85$, 1000 training and 100 test samples, 500 epochs, one training run each): mean (median) relative $L^2$ test error in per cent, the value published for the same architecture, and their ratio. Implementations: FNO-port, faithful port of \cite{li2021fno}; FNO-var, the simplified variant behind our FNO-op rows; MgNO-off, the official MgNO model classes \cite{he2024mgno}; Transolver-re, our re-implementation of \cite{wu2024transolver}; WNO-db4 and WNO-db6, ports of version 2.0.0 of the authors' WNO code with the db4 transform of the paper and the db6 transform of the repository; WNO-light, our light WNO layer. Published values: FNO \cite{li2021fno}; WNO, db4 with four levels \cite{tripura2023wno}; Transolver \cite{wu2024transolver}, obtained with 200 test samples and an additional gradient term in the loss. The MgNO paper \cite{he2024mgno} reports a different smooth Darcy data set, so no value is listed. ``Used in'': the rows of Tables~\ref{tab:channels} and~\ref{tab:barriers} (or the excluded run of Table~\ref{tab:excluded}) that use the same implementation.}\label{tab:published}
\begin{adjustbox}{max width=\linewidth}
\begin{tabular}{llrlll}\toprule
implementation & used in & params & rel.\ $L^2$ (\%) & published (\%) & ratio \\\midrule
FNO-port & -- & 1{,}188{,}353 & 0.96 (0.88) & 1.08 & 0.88 \\
FNO-var & FNO-op & 595{,}265 & 1.74 (1.60) & 1.08 & 1.61 \\
MgNO-off & MgNO, MgNO-op & 645{,}828 & 0.70 (0.62) & -- & -- \\
Transolver-re & Transolver-op & 976{,}961 & 1.22 (1.08) & 0.57 & 2.13 \\
WNO-db4 & -- & 7{,}955{,}201 & 2.27 (2.10) & 0.84 & 2.70 \\
WNO-db6 & excl. & 14{,}770{,}945 & 2.14 (2.03) & 0.84 & 2.55 \\
WNO-light & WNO-op & 334{,}209 & 3.53 (3.31) & -- & -- \\
\bottomrule\end{tabular}
\end{adjustbox}
\end{table}

The faithful FNO port attains an error comparable to the published benchmark value ($0.96\%$ against $1.08\%$), and the
official MgNO classes train successfully on this data set ($0.70\%$; the MgNO paper reports a different smooth Darcy data set, so
no published value is compared). These are useful checks of those implementations and of the benchmark pipeline; they do not
validate every adaptation used in the unstructured-mesh experiments, where the mesh--grid bridge (\appref{app:bridge}), the
energy-containing loss, the matched parameter budget and the contrast all differ, and the faithful port is not the FNO
implementation of the FNO-op rows. That variant reaches $1.74\%$ with $595{,}265$ parameters against the port's $1{,}188{,}353$;
the comparison measures the aggregate difference between these two configurations in this run, not the effect of any single
simplification. Three implementations do not reach the published accuracy. Our Transolver re-implementation reaches $1.22\%$
against the published $0.57\%$, which was, however, obtained with an additional gradient term in the loss and 200 test samples.
The port of the published WNO reaches $2.27\%$ with db4 and $2.14\%$ with db6 against $0.84\%$, although its final training loss
with db4 ($0.80\%$) is close to the published test error; we did not establish the cause of this gap. Our light WNO layer reaches
$3.53\%$. No published-benchmark check was run for LiNO. The Transolver-op, WNO-op and LiNO-op rows of the main tables therefore
measure these implementations only.

\section{The grid bridge}\label{app:bridge}
\begin{table}[tbp]\centering\footnotesize\setlength{\tabcolsep}{4pt}
\caption{The grid bridge on 40 training operators (channels, seed 1): numerical rank of the grid-to-mesh map $T$ (eigenvalues of $TT^\top$ above $10^{-12}$ times the largest); energy error of the round trip of the exact correction $t=u-u_c$ (mesh $\to$ grid by piecewise-linear interpolation, grid $\to$ mesh by bilinear interpolation), relative to $\|t\|_A$ and to $\|u\|_A$; and of the best approximation of $t$ in the numerically resolved range of $T$ (Galerkin solve on an orthonormal basis of that range). Mean relative energy errors over the samples.}\label{tab:bridge}
\begin{adjustbox}{max width=\linewidth}
\begin{tabular}{llllll}\toprule
grid & rank($T$) & round trip / $\|t\|_A$ & round trip / $\|u\|_A$ & best / $\|t\|_A$ & best / $\|u\|_A$ \\\midrule
$63\times63$ & 3126 of 3486 & 12.62 & 4.77 & 3.25e-01 & 1.23e-01 \\
$90\times90$ & 3486 of 3486 & 9.26 & 3.50 & 5.45e-12 & 2.07e-12 \\
$126\times126$ & 3486 of 3486 & 6.79 & 2.56 & 2.54e-13 & 9.66e-14 \\
\bottomrule\end{tabular}
\end{adjustbox}
\end{table}

On 40 training operators, sampling the exact correction $t=u-u_c$ on the protocol's $63\times63$ grid and interpolating it back
gives a mean relative energy error of $12.6$ with respect to $t$, with a small nodal error. This round-trip error is not a lower
bound. The range of the grid-to-mesh map $T$ is resolved numerically: its basis consists of the eigenvectors of $TT^\top$ (on the
interior nodes) with eigenvalues above $10^{-12}$ times the largest. At the protocol's grid the map is rank deficient (numerical
rank 3126 of 3486; the mesh is locally finer than the grid near the holes). For a fixed sample, the minimum over all grid outputs
$g$ of $\|t-Tg\|_A/\|t\|_A$ is a lower bound on the error of any prediction restricted to that range; the $A$-orthogonal
projection onto the numerically resolved range estimates it, and on these training samples its mean is $0.325$ with respect to
the correction, or $0.123$ with respect to the full solution. The estimate is not certified for the full algebraic range, and the
sample averages are not guarantees for unseen operators. On finer grids the map has full numerical row rank, so any nodal field is
representable in principle, but the round trip of the exact correction, the field a smooth grid model would naturally produce,
still has mean relative energy errors several times one (Table~\ref{tab:bridge}).

\section{Repeated training and the output/loss ablation}\label{app:repeat}
\begin{table}[tbp]\centering\footnotesize\setlength{\tabcolsep}{4pt}
\caption{Repeated training with initialisation seeds 0 (the runs of the other tables), 1 and 2 on the same data. (a) OA-VCB: reference/OA-VCB $G$ ($W$) per data seed and initialisation seed, and the range of the mean energy error over the three trained models. (b) The residual-linear control OA-VCBL on channels, data seed 1: mean energy error of its projected (deployed) and plain cycle output, reference/OA-VCBL $G$ ($W$), and OA-VCBL/OA-VCB $G$ ($W$) with OA-VCB of the same initialisation seed (the two factors of Section~\ref{sec:ablation}). The ranges describe variability across training runs; the three models are evaluated on the same 100 test operators, so they are not independent test observations.}\label{tab:initseeds}
{\small (a) OA-VCB}\par\smallskip
\begin{adjustbox}{max width=\linewidth}
\begin{tabular}{llllll}\toprule
class & data seed & init.\ 0 & init.\ 1 & init.\ 2 & mean energy (range) \\\midrule
channels & 1 & 10.23 (1.00) & 10.38 (1.00) & 10.74 (1.00) & 2.43e-3--2.55e-3 \\
channels & 2 & 9.76 (1.00) & 9.94 (1.00) & 10.08 (1.00) & 2.72e-3--2.79e-3 \\
channels & 3 & 9.68 (1.00) & 9.96 (1.00) & 10.30 (1.00) & 2.79e-3--2.92e-3 \\
\midrule
barriers & 1 & 2.13 (0.92) & 2.09 (0.94) & 2.04 (0.91) & 4.36e-2--4.61e-2 \\
barriers & 2 & 1.93 (0.83) & 1.93 (0.85) & 1.96 (0.85) & 6.64e-2--6.96e-2 \\
barriers & 3 & 1.93 (0.89) & 1.92 (0.90) & 1.80 (0.86) & 6.06e-2--6.30e-2 \\
\bottomrule\end{tabular}
\end{adjustbox}
\par\medskip
{\small (b) OA-VCBL, channels, data seed 1}\par\smallskip
\begin{adjustbox}{max width=\linewidth}
\begin{tabular}{lllll}\toprule
init.\ seed & projected & plain & ref./OA-VCBL & OA-VCBL/OA-VCB \\\midrule
0 & 1.01e-2 & 1.07e-2 & 2.72 (0.99) & 3.76 (1.00) \\
1 & 1.01e-2 & 1.07e-2 & 2.72 (0.99) & 3.81 (1.00) \\
2 & 9.83e-3 & 5.87e-2 & 2.63 (1.00) & 4.08 (1.00) \\
\bottomrule\end{tabular}
\end{adjustbox}
\end{table}

\begin{table}[tbp]\centering\footnotesize\setlength{\tabcolsep}{4pt}
\caption{Plain and projected outputs of the same trained weights (mean energy error, held-out seed 1). Models with loss ``plain'' optimise the plain output only; models with loss ``plain + proj.'' optimise both output losses, and their projected output is the one deployed. In each row the other column is the same network evaluated with the other output.}\label{tab:factorial8}
\begin{adjustbox}{max width=\linewidth}
\begin{tabular}{lllllll}\toprule
class & model & deflation & loss & smoothers & plain output & projected output \\\midrule
barriers & OA-VC & no & plain & learned & 0.2329 & 0.1206 \\
channels & OA-VC & no & plain & learned & 0.0041 & 0.0035 \\
barriers & OA-VCZ & yes & plain & learned & 0.0751 & 0.0487 \\
channels & OA-VCZ & yes & plain & learned & 0.0035 & 0.0033 \\
barriers & OA-VCQ & no & plain + proj. & learned & 0.3111 & 0.1072 \\
barriers & OA-VCB & yes & plain + proj. & learned & 0.0663 & 0.0436 \\
channels & OA-VCB & yes & plain + proj. & learned & 0.0026 & 0.0026 \\
barriers & OA-VCL & yes & plain & residual-linear & 0.1508 & 0.0683 \\
channels & OA-VCL & yes & plain & residual-linear & 0.0111 & 0.0104 \\
barriers & OA-VCBL & yes & plain + proj. & residual-linear & 0.1518 & 0.0676 \\
channels & OA-VCBL & yes & plain + proj. & residual-linear & 0.0107 & 0.0101 \\
\bottomrule\end{tabular}
\end{adjustbox}
\end{table}

Table~\ref{tab:initseeds} reports the three initialisation seeds of OA-VCB on every held-out data set and of the residual-linear
control OA-VCBL on channels (data seed 1). One OA-VCBL run is not a repeat of the others in the optimisation sense: initialisation
seed 1 reaches almost the same final training loss as seed 0, while seed 2 ends with a $2.8$ times higher final training loss and a
much larger plain-cycle error, with a similar projected mean error and a slightly lower paired gain (Section~\ref{sec:ablation}).
Matching initialisation seed numbers across architectures index runs; they do not match optimisation trajectories.
Table~\ref{tab:factorial8} evaluates each trained network of Table~\ref{tab:factorial} with both outputs, which separates the
output used at inference from the training objective. Projecting the output of a network trained without the projection term
also helps on its own (OA-VC $0.233\to0.121$, OA-VCZ $0.075\to0.049$ on barriers). One channel sample (id 320) is difficult for
every model without the trained Galerkin output and for OA-VCBL; we do not interpret it further.

\section{Sensitivity of the detector thresholds and of the transfer steps}\label{app:sensitivity}
Table~\ref{tab:sensitivity} varies two choices that were fixed during development, using the classical reference, which needs
no retraining (\path{diagnostics_cost/sensitivity_detector_transfers.py}). These are subsequent checks; the main experiments keep
the development values $\theta = 0.1$, $\tau = 0.05$. Recall counts islands (202 in the 300 held-out operators; an operator may
contain several) whose indicator lies in $\operatorname{span}Z$ to a relative Euclidean distance below $10^{-8}$. Precision counts
detected vectors with at least $90\%$ of their support inside one island, out of all detected vectors; there is no one-to-one
matching, so the two numerators count different objects. (a) The threshold $\theta$ matters more than $\tau$. At the default
$\tau = 0.05$, $\theta = 0.05$, $0.1$ and $0.2$ recover $202$, $195$ and $81$ of the $202$ islands; $\theta = 0.05$ does so with one
detected vector outside an island, and $\theta = 0.2$ more than doubles the reference's error on island samples. Varying $\tau$
in $[0.02, 0.1]$ changes the reference's mean error less than varying $\theta$, but it can change the detected space: at
$\theta = 0.1$ the recall rises from $195$ to $201$ and the number of detected vectors from $204$ to $224$, and at $\theta = 0.2$
the number of detected vectors grows from $116$ to $498$ while the error changes little. Only at $\theta = 0.05$ are the tested
values of $\tau$ without effect. (b) Re-evaluating the trained barrier models with the $Z$ of $\theta = 0.05$ (only deflation
and the Galerkin output change) lowers the mean error of OA-VCB by $4.6\%$ and that of the reference by $3.1\%$; the pooled
paired gain changes from $2.00$ to $2.03$ (\path{diagnostics_cost/detector_theta_learned.py}). (c) More Jacobi steps in the
transfers improve the classical reference, by $7\%$ on channels and $30\%$ on barriers from $m = 4$ to $m = 8$, while the
operator complexity grows from $4.5$ to $7.4$ and the set-up time fourfold; with $m = 1$ or $2$ the reference loses $7$--$11\%$ on
channels and $38$--$77\%$ on barriers. Retraining OA-VCB with $m = 8$ (held-out seed 1, Table~\ref{tab:m8}) keeps the learned
margin over the correspondingly stronger reference: $G = 11.17$ on channels ($10.23$ with $m = 4$) and $2.01$ on barriers
($2.13$), with $W = 1.00$ and $0.90$. Measured against the $m = 4$ reference, the $m = 8$ model gains $12.0$ and $2.88$.
\begin{table}[tbp]\centering\footnotesize\setlength{\tabcolsep}{4pt}
\caption{Transfers with $m=8$ Jacobi steps (held-out data seed 1, initialisation seed 0): OA-VCB trained with $m=8$ and its classical reference with $m=8$; mean energy errors, paired $G$ [95\% CI] and $W$; for comparison the $m=4$ margin and the improvement of the classical reference from $m=4$ to $m=8$.}\label{tab:m8}
\begin{adjustbox}{max width=\linewidth}
\begin{tabular}{lllllll}\toprule
class & OA-VCB ($m=8$) & reference ($m=8$) & $G$ [95\% CI] & $W$ & $G$ at $m=4$ & ref.\ $m=4$/$m=8$ \\\midrule
channels & 2.14e-3 & 2.42e-2 & 11.17 [10.74, 11.63] & 1.00 & 10.23 [9.75, 10.71] & 1.07 \\
barriers & 3.48e-2 & 5.19e-2 & 2.01 [1.79, 2.23] & 0.90 & 2.13 [1.90, 2.38] & 1.44 \\
\bottomrule\end{tabular}
\end{adjustbox}
\end{table}

\begin{table}[tbp]\centering\footnotesize\setlength{\tabcolsep}{4pt}
\caption{Sensitivity of two fixed choices (classical reference unless stated; L3, $k_0=2$). (a) Detector thresholds on the barrier test sets of held-out seeds 1--3 (300 operators, 202 coefficient islands in total; an operator may contain several). Recall: islands whose indicator lies in $\operatorname{span}Z$ (relative Euclidean distance below $10^{-8}$), out of all islands. Precision: detected vectors with at least $90\%$ of their support inside one island, out of all detected vectors (no one-to-one matching, so the two numerators count different objects and need not agree). Mean reference energy error over all samples and over the island samples. (b) The trained barrier OA-VCB models (trained with $\theta=0.1$) evaluated with $Z$ from $\theta=0.1$ and $0.05$ ($\tau=0.05$), pooled seeds 1--3: mean energy errors, reference/OA-VCB $G$ [95\% CI] and $W$. (c) Number $m$ of Jacobi steps in the transfers, held-out seed 1: operator complexity $\sum_k\operatorname{nnz}(A_k)/\operatorname{nnz}(A_L)$ (median), hierarchy build time (median, one CPU), reference energy error and the paired geometric-mean ratio of its errors to those at $m=4$; the learned cycle was trained only with $m=4$ and $m=0$ (Table~\ref{tab:transfers}).}\label{tab:sensitivity}
{\small (a) detector thresholds, classical reference}\par\smallskip
\begin{adjustbox}{max width=\linewidth}
\begin{tabular}{llllll}\toprule
$\theta$ & $\tau$ & recall & precision & energy, all & energy, island samples \\\midrule
0.05 & 0.02 & 202/202 & 202/203 & 7.86e-2 & 1.65e-1 \\
0.05 & 0.05 & 202/202 & 202/203 & 7.86e-2 & 1.65e-1 \\
0.05 & 0.1 & 202/202 & 202/203 & 7.86e-2 & 1.65e-1 \\
0.1 & 0.02 & 195/202 & 197/204 & 8.11e-2 & 1.72e-1 \\
0.1 & 0.05 (used) & 195/202 & 199/211 & 8.11e-2 & 1.72e-1 \\
0.1 & 0.1 & 201/202 & 208/224 & 7.86e-2 & 1.66e-1 \\
0.2 & 0.02 & 79/202 & 95/116 & 1.97e-1 & 4.43e-1 \\
0.2 & 0.05 & 81/202 & 131/211 & 1.95e-1 & 4.39e-1 \\
0.2 & 0.1 & 86/202 & 187/498 & 1.91e-1 & 4.30e-1 \\
\bottomrule\end{tabular}
\end{adjustbox}
\par\medskip
{\small (b) trained OA-VCB with the detector of a smaller $\theta$}\par\smallskip
\begin{adjustbox}{max width=\linewidth}
\begin{tabular}{lllll}\toprule
$\theta$ & OA-VCB & reference & $G$ [95\% CI] & $W$ \\\midrule
0.1 & 5.80e-2 & 8.11e-2 & 2.00 [1.86, 2.14] & 0.88 \\
0.05 & 5.53e-2 & 7.86e-2 & 2.03 [1.90, 2.17] & 0.89 \\
\bottomrule\end{tabular}
\end{adjustbox}
\par\medskip
{\small (c) Jacobi steps in the transfers, classical reference}\par\smallskip
\begin{adjustbox}{max width=\linewidth}
\begin{tabular}{lllllll}\toprule
$m$ & op.\ complexity & build (ms) & channels & $G$ vs $m=4$ & barriers & $G$ vs $m=4$ \\\midrule
0 & 1.27 & 2 & 3.23e-2 & 1.24 & 1.23e-1 & 2.82 \\
1 & 1.82 & 6 & 2.88e-2 & 1.11 & 1.02e-1 & 1.77 \\
2 & 2.63 & 13 & 2.78e-2 & 1.07 & 9.23e-2 & 1.38 \\
4 (used) & 4.54 & 56 & 2.59e-2 & 1.00 & 7.51e-2 & 1.00 \\
8 & 7.44 & 232 & 2.42e-2 & 0.93 & 5.19e-2 & 0.70 \\
\bottomrule\end{tabular}
\end{adjustbox}
\end{table}


\section{Harmonic enrichment: protocol, amendment and results}\label{app:harmonic}
\begin{table}[tbp]\centering\footnotesize\setlength{\tabcolsep}{4pt}
\caption{Post-hoc harmonic enrichment of the frozen OA-VCB predictions (barriers; patch radius 2, fixed on the development set; the fresh confirmation set has 100 new operators per checkpoint). (a) Detector-positive samples / all samples; mean energy error before and after the local correction; $R = 1-\overline{e}_{\mathrm{after}}/\overline{e}_{\mathrm{before}}$ in \% [95\% CI]; $G_{24}$: paired ratio of the patch-indicator control's error to the harmonic correction's [95\% CI]. (b) Per-sample captured fraction of squared energy error, median [quartiles]; on the fixed global-reference subset (number of samples), $R$ of the local and of the global correction, in \%.}\label{tab:harmonic}
{\small (a) Primary outcome}\par\smallskip
\begin{adjustbox}{max width=\linewidth}
\begin{tabular}{llllll}\toprule
set & samples & $e$ before & $e$ after & $R$ (\%) & $G_{24}$ \\\midrule
development (seed 1) & 42/100 & 0.0956 & 0.0876 & 8.3 [5.1, 12.3] & 1.10 [1.06, 1.13] \\
retrospective (seeds 2, 3) & 92/200 & 0.1352 & 0.1170 & 13.5 [9.6, 17.9] & 1.15 [1.11, 1.19] \\
fresh confirmation & 89/200 & 0.1219 & 0.1043 & 14.4 [9.8, 19.5] & 1.13 [1.09, 1.16] \\
\bottomrule\end{tabular}
\end{adjustbox}
\par\medskip
{\small (b) Captured fraction and local versus global extension}\par\smallskip
\begin{adjustbox}{max width=\linewidth}
\begin{tabular}{lllll}\toprule
set & captured fraction & subset size & local $R$ (\%) & global $R$ (\%) \\\midrule
development (seed 1) & 0.09 [0.03, 0.31] & 20 & 9.0 & 17.6 \\
retrospective (seeds 2, 3) & 0.13 [0.03, 0.40] & 40 & 12.2 & 16.2 \\
fresh confirmation & 0.12 [0.02, 0.33] & 40 & 10.9 & 20.9 \\
\bottomrule\end{tabular}
\end{adjustbox}
\end{table}

\paragraph{Correction.} For a frozen prediction $x$ with residual $r = F - Ax$ and a small matrix $Z$ built from $A$ alone
(independent columns), the correction
\begin{gather*}
  x_+ = x + Z(Z^\top AZ)^{-1}Z^\top r,\\
  \|u-x_+\|_A^2 = \|u-x\|_A^2 - (Z^\top r)^\top(Z^\top AZ)^{-1}(Z^\top r),
\end{gather*}
is the $A$-orthogonal projection of the error onto $\operatorname{span}Z$ (exact arithmetic): a sequential correction of the
frozen error, not a joint re-optimisation over the cycle's own directions together with $Z$; the networks are not retrained. An
earlier classical pilot had suggested that an $A$-harmonic tail approximates the variation of the slowest mode across the
barrier; a low eigenmode need not satisfy $Az=0$ outside the island, but the enrichment constructed here is harmonic by
definition. $Z$ contains the indicators and, for each island, its $A$-harmonic extension into a patch of $\rho$ graph rings (zero
outside). A control of equal support replaces the harmonic extension by the indicator of the patch, and a reference uses the
global harmonic extension. Island samples here are those on which the detector of Section~\ref{sec:deflation} returns at least
one vector (42, 92 and 89 samples in the three sets, against 39, 88 and 84 with the geometric label). We report
$R = 1-\overline{e}_{\mathrm{after}}/\overline{e}_{\mathrm{before}}$, the reduction of the mean relative energy error; $G_{24}$,
the paired ratio of the control's error to the harmonic correction's; and the per-sample captured fraction of squared energy
error, $\kappa = 1 - \|u-x_+\|_A^2/\|u-x\|_A^2$.

\paragraph{Protocol.} The protocol (version 2) fixed, before any enrichment number: the frozen checkpoints and the inference
path; the three data sets and their roles (development, retrospective, fresh confirmation); the island label (the detector of
Section~\ref{sec:deflation} with $\theta=0.1$, $\tau=0.05$; an independent geometric label reported only); the corrections
(indicators; local harmonic extensions for $\rho\in\{1,2,4,8\}$; the patch-indicator control; global extensions on the first 20
detector-positive samples per checkpoint, i.e.\ per data seed: 20 / 40 / 40 samples in the three sets); the numerics (columns
normalised in the energy norm, greedy pivoted Cholesky of the scaled Gram matrix with pivot threshold $10^{-10}$; gain computed
as $2d^\top r - d^\top Ad$ and checked against the change of squared error, agreement to $10^{-10}$ relative); the radius rule
($\rho = 2$, the smallest radius within $10\%$ of the best on the seed-1 test set); and the outcome thresholds, ``strong support''
($R\ge0.5$ and $G_{24}\ge1.25$), ``partial support'' ($R\ge0.2$) and ``little support'' otherwise, decided on a fresh
confirmation set of $100$ new coefficient fields per checkpoint, never generated before. Seeds 2 and 3 served as a retrospective
evaluation: they were held out from training and from the radius choice, but their model errors had been inspected before.
\emph{Amendment (version 3), made after the reproduction check and before any enrichment number.} The planned check required the
CPU re-evaluation of the stored predictions to agree with the GPU records to $10^{-4}$ per sample. Means agreed to six digits on
all three seeds and samples without islands to $3\times10^{-5}$, but island samples differed by up to $1.4\times10^{-3}$. On the
same CPU, changing only the batch size moved island samples by up to $1.1\times10^{-4}$ and evaluating the network in double
precision by up to $1.75\times10^{-4}$, consistent with single-precision round-off amplified by near-singular operators, although
these checks do not show that round-off accounts for the whole CPU/GPU difference. The study therefore used the CPU
single-precision path; the criterion was replaced by agreement of the means to $10^{-5}$ and of the non-island samples to $10^{-4}$.

\paragraph{Further results.} The captured fraction is heterogeneous (median $0.12$, quartiles $0.02$ and $0.33$ on the fresh
operators). The classical reference behaves similarly ($R = 18.5\%$ locally on the fresh operators; $15.6\%$ local and $29.2\%$
global on the reference subset; per-seed values, including one seed with partial support, in Table~\ref{tab:harmonicseeds}). The
indicator columns alone gave numerically negligible gains (at most $10^{-19}$ relative), consistent with both outputs already
containing their span. As an observation on the 42 detector-positive development samples, the means of the per-sample regional
fractions of squared energy error are $28\%$ inside the islands, $31\%$ in the barriers and $41\%$ in the connected conductive
region; after the correction the barrier fraction drops to $20\%$. Building the local basis took about $0.6$\,ms and applying the
projection, including the study's diagnostics, about $0.5$\,ms per island sample (median; the detection, $1.5$\,ms, is already
part of OA-VCB).
\begin{table}[tbp]\centering\footnotesize\setlength{\tabcolsep}{4pt}
\caption{Harmonic enrichment per seed, OA-VCB and classical reference: detector-positive samples, $R$ in \% [95\% CI], $G_{24}$, and the outcome category of the protocol.}\label{tab:harmonicseeds}
\begin{adjustbox}{max width=\linewidth}
\begin{tabular}{lllllll}\toprule
set & seed & predictor & samples & $R$ (\%) & $G_{24}$ & outcome \\\midrule
development & 1 & OA-VCB & 42 & 8.3 [5.1, 12.3] & 1.10 & little \\
development & 1 & reference & 42 & 18.9 [14.8, 23.7] & 1.24 & little \\
retrospective & 2 & OA-VCB & 49 & 17.9 [11.6, 24.5] & 1.19 & little \\
retrospective & 2 & reference & 49 & 16.4 [12.2, 21.0] & 1.22 & little \\
retrospective & 3 & OA-VCB & 43 & 8.4 [5.3, 11.8] & 1.11 & little \\
retrospective & 3 & reference & 43 & 14.3 [9.9, 19.4] & 1.21 & little \\
fresh confirmation & 2 & OA-VCB & 37 & 17.8 [9.8, 26.1] & 1.15 & little \\
fresh confirmation & 2 & reference & 37 & 22.7 [18.0, 27.8] & 1.31 & partial \\
fresh confirmation & 3 & OA-VCB & 52 & 11.5 [6.9, 16.6] & 1.11 & little \\
fresh confirmation & 3 & reference & 52 & 15.3 [10.6, 20.8] & 1.22 & little \\
\bottomrule\end{tabular}
\end{adjustbox}
\end{table}


\section{Load tests}\label{app:loads}
\paragraph{Loads.} The random-load distribution used for training (Section~\ref{sec:k0}) draws a constant $c\sim U(-1,1)$ plus a
sine expansion $\sum_{p,q=0}^{3}K_{pq}\sin((p+1)\pi x)\sin((q+1)\pi y)/(1+p+q)$ with $K_{pq}\sim N(0,1)$, to which, with
probability one half, a Gaussian of width $U(0.02,0.08)$ and amplitude $U(2,10)$ centred in $[0.1,0.9]^2$ is added. The smooth
test uses the same finite sine span with a different coefficient distribution and no constant term. The localised test is a
pure, unit-amplitude Gaussian source of fixed width $0.03$ centred in $[0.15,0.85]^2$; its width lies inside the training
interval, so it changes the composition of the forcing rather than its width, and by positive homogeneity its amplitude is
immaterial. The load values $f=\pm1$ lie inside the support of the training distribution, so they test coverage of that
distribution's extreme members; neither test is extrapolation beyond the support, and a different generating procedure does not
by itself place a load outside it. A miss of the per-seed criterion by one sample is a threshold outcome, not evidence of a sharp
boundary.
\begin{table}[tbp]\centering\footnotesize\setlength{\tabcolsep}{4pt}
\caption{The fixed-load ($f=1$) OA-VCB checkpoints of Table~\ref{tab:core} applied to other loads (channels, held-out seeds 1--3, 300 operators); $u_c$ is recomputed for each load. Mean energy errors, reference/OA-VCB $G$ per seed and pooled, $W$, and the number of samples on which OA-VCB is worse than $u_c$.}\label{tab:rhs}
\begin{adjustbox}{max width=\linewidth}
\begin{tabular}{llllllll}\toprule
load & $u_c$ & OA-VCB & reference & $G$ per seed & $G$ pooled [95\% CI] & $W$ & $>u_c$ \\\midrule
$f=1$ (training load) & 0.373 & 2.75e-3 & 2.57e-2 & 10.23 / 9.76 / 9.68 & 9.89 [9.57, 10.20] & 1.00 & 0 \\
$f=-1$ & 0.373 & 1.37e-1 & 2.57e-2 & 0.30 / 0.28 / 0.11 & 0.21 [0.20, 0.23] & 0.00 & 0 \\
$f=3$ & 0.373 & 2.75e-3 & 2.57e-2 & 10.23 / 9.76 / 9.68 & 9.89 [9.57, 10.20] & 1.00 & 0 \\
random-load distribution (Sec.~\ref{sec:k0}) & 0.397 & 1.01e-1 & 2.59e-2 & 0.66 / 0.74 / 0.45 & 0.60 [0.50, 0.72] & 0.37 & 0 \\
smooth, both signs & 0.386 & 1.55e-1 & 2.49e-2 & 0.25 / 0.29 / 0.12 & 0.20 [0.19, 0.23] & 0.06 & 0 \\
localised source & 0.635 & 1.08e-1 & 2.21e-2 & 0.38 / 0.54 / 0.42 & 0.44 [0.38, 0.51] & 0.32 & 0 \\
\bottomrule\end{tabular}
\end{adjustbox}
\end{table}

\begin{table}[tbp]\centering\footnotesize\setlength{\tabcolsep}{4pt}
\caption{The OA-VCB checkpoints trained on the random-load distribution (Section~\ref{sec:k0}; one model per held-out data seed 1--3) on the loads of Table~\ref{tab:rhs}, same operators. (a) Pooled over the 300 operators: mean energy errors of the random-load models and of the reference, reference/random-model $G$ [95\% CI] and $W$, and samples worse than $u_c$. (b) Per seed (1 / 2 / 3): reference/random-model $G$ and $W$; mean error of the fixed-load models of Table~\ref{tab:rhs}; and pooled fixed-load/random-load $G$ and $W$ (above 1 and above 0.5 favour the random-load models). The prespecified criterion per seed and load family was $G\ge2$ and $W\ge0.8$.}\label{tab:rhsrandom}
{\small (a) Pooled}\par\smallskip
\begin{adjustbox}{max width=\linewidth}
\begin{tabular}{llllll}\toprule
load & random-load & reference & $G$ [95\% CI] & $W$ & $>u_c$ \\\midrule
$f=1$ & 3.99e-3 & 2.57e-2 & 6.87 [6.62, 7.13] & 1.00 & 0 \\
$f=-1$ & 4.21e-3 & 2.57e-2 & 6.51 [6.27, 6.75] & 1.00 & 0 \\
$f=3$ (homogeneity check) & 3.99e-3 & 2.57e-2 & 6.87 [6.62, 7.13] & 1.00 & 0 \\
random-load distribution (Sec.~\ref{sec:k0}) & 4.43e-3 & 2.59e-2 & 6.22 [5.99, 6.44] & 1.00 & 0 \\
smooth, both signs & 4.00e-3 & 2.49e-2 & 6.38 [6.20, 6.56] & 1.00 & 0 \\
localised source & 1.25e-2 & 2.21e-2 & 2.51 [2.25, 2.79] & 0.84 & 0 \\
\bottomrule\end{tabular}
\end{adjustbox}
\par\medskip
{\small (b) Per seed, and against the fixed-load models}\par\smallskip
\begin{adjustbox}{max width=\linewidth}
\begin{tabular}{llllll}\toprule
load & $G$ per seed & $W$ per seed & fixed-load & fixed/random $G$ & $W$ \\\midrule
$f=1$ & 6.90 / 6.75 / 6.98 & 1.00 / 1.00 / 1.00 & 2.75e-3 & 0.70 & 0.01 \\
$f=-1$ & 6.67 / 6.39 / 6.47 & 1.00 / 1.00 / 1.00 & 1.37e-1 & 30.8 & 1.00 \\
$f=3$ (homogeneity check) & 6.90 / 6.75 / 6.98 & 1.00 / 1.00 / 1.00 & 2.75e-3 & 0.70 & 0.01 \\
random-load distribution (Sec.~\ref{sec:k0}) & 6.27 / 6.05 / 6.34 & 1.00 / 1.00 / 1.00 & 1.01e-1 & 10.3 & 0.86 \\
smooth, both signs & 6.41 / 6.31 / 6.42 & 1.00 / 1.00 / 1.00 & 1.55e-1 & 31.1 & 1.00 \\
localised source & 2.80 / 2.15 / 2.61 & 0.89 / 0.79 / 0.83 & 1.08e-1 & 5.7 & 0.99 \\
\bottomrule\end{tabular}
\end{adjustbox}
\end{table}


\section{Further settings}\label{app:settings}
\paragraph{No external coarse solution ($k_0 = 0$).} The networks start from $u=0$; the network and the reference still contain
their exact 28-node coarsest solve. Two-grid PCG from zero first reaches OA-VCB's error after a median of $5$ iterations. The
operator-informed baselines are again at or below the reference: OA-Lift at its level (reference/OA-Lift $G = 1.02$--$1.05$) and
GNN-op about half as accurate.
\paragraph{Finer hierarchy.} $L=4$, $k_0=3$, 14{,}398 unknowns, one held-out seed; two-grid PCG first reaches OA-VCB's error after a
median of $6$ iterations (range $4$--$8$). This is not a mesh-refinement study at fixed information ($k_0$ changes with $L$). The
L4 OA-Lift run had only 13{,}215 parameters and is not reported as a matched baseline.
\paragraph{Refinement without retraining.} The L4-trained model was applied unchanged to the same coefficient draws on the next
two refined meshes of held-out seed 1, with the reference on the same hierarchy. At L5 (58{,}494 unknowns, $k_0 = 4$; the
four-level cycle now ends on mesh level 1 instead of 0) it keeps $G = 8.12$ [7.50, 8.67] with $W = 1$ on all 100 test samples
(smallest per-sample ratio $1.65$) and is never worse than $u_c$; paired with the L4 results for the same coefficient draws
($G = 9.72$), its per-sample gain is $0.84$ [0.76, 0.91] times the L4 gain, smaller at L5 on $81\%$ of the samples. At L6
(235{,}774 unknowns, $k_0 = 5$, bottom on mesh level 2) the evaluation was stopped by the session time limit of the computing
platform after the first 21 of the 24 planned test samples (ids 300--320) and wrote no per-sample record; the values are taken from
the per-sample lines of the run log (four significant digits; rounding moves $G$ by less than $0.005$). On these completed cases
$G = 4.18$ [3.43, 4.88] with $W = 1$ and a smallest per-sample ratio of $1.06$, about half the L5 gain on the same draws (paired
ratio $0.51$ [0.46, 0.56]). The interval describes the completed cases; it does not account for the missing cases. The original
logs and the L5 record are kept with the code archive (\path{Results_kaggle_session 11/}). A model trained at L3 does not keep its margin
at L4 ($G = 1.95$, $W = 0.95$, 20 samples; $1.97$ with the cycle depth matched). In this evaluation, operator data and inference took
about $64$\,s per sample at L5 and $1{,}370$\,s at L6 on one CPU shared by concurrent jobs; these are not comparable with
Table~\ref{tab:cost}.
\paragraph{Polygonal domain.} A non-convex polygon with a re-entrant corner, a different unstructured coarse mesh with a 45-node
coarsest level: $G = 9.39$, $W=1$, one held-out seed.
\paragraph{Random loads, in distribution.} Trained and tested on the distribution of \appref{app:loads}, one training run per data
seed: $G = 6.27$ [5.89, 6.65] on seed 1, which meets the outcome rule fixed before the run (``strong'' if $G\ge2$ and $W\ge0.8$), and
$6.05$ and $6.34$ on seeds 2 and 3, with the same rule fixed beforehand for both, all with $W = 1$.
\paragraph{Anisotropic tensor coefficient.} Meandering fibre tows with $K = k_T I + (k_L - k_T)\,t\,t^\top$ along the band direction
$t$, $k_L = 100$, $k_T = 5$, unit background, insulated holes: $G = 3.51$ and $3.58$ on held-out seeds $1$, $2$, with $W=1$ (seed $0$,
$3.50$, is reported although OA-VCB was never developed on this class). Here the reference is already more accurate than the plain
classical cycle (ratio of mean errors $1.37$). The detector returns at least one vector on $9$--$10\%$ of the samples (at most two);
with insulated holes the coarsest level has 52 nodes.
\paragraph{Plane-strain elasticity (preliminary).}\label{sec:elasticity}
\begin{table}[tbp]\centering\footnotesize\setlength{\tabcolsep}{4pt}
\caption{Plane-strain elasticity (stiff bands, $E_f/E_m=100$, traction-free pores; one data set, test samples 300--399; preliminary). Mean energy error, and paired reference/model $G$ [95\% CI] and $W$. OA-VCB has one training run; the two baselines were trained with three initialisation seeds on the same data. Two-grid PCG first reaches OA-VCB's per-sample energy error after a median of 11 iterations (range 4--32).}\label{tab:elasticity}
\begin{adjustbox}{max width=\linewidth}
\begin{tabular}{lrlll}\toprule
model & params & energy & $G$ [95\% CI] & $W$ \\\midrule
OA-VCB (ours) & 20{,}672 & 2.47e-2 & 2.32 [2.20, 2.44] & 1.00 \\
OA-Lift, init.\ seed 0 & 21{,}910 & 2.11e-1 & 0.25 [0.24, 0.26] & 0.00 \\
OA-Lift, init.\ seed 1 & 21{,}910 & 2.18e-1 & 0.24 [0.23, 0.26] & 0.00 \\
OA-Lift, init.\ seed 2 & 21{,}910 & 2.12e-1 & 0.25 [0.23, 0.26] & 0.00 \\
GNN-op, init.\ seed 0 & 20{,}354 & 1.96e-1 & 0.27 [0.25, 0.28] & 0.00 \\
GNN-op, init.\ seed 1 & 20{,}354 & 1.95e-1 & 0.27 [0.25, 0.28] & 0.00 \\
GNN-op, init.\ seed 2 & 20{,}354 & 1.96e-1 & 0.27 [0.25, 0.28] & 0.00 \\
\midrule
classical reference & -- & 5.50e-2 & 1 & -- \\
classical V(4,4), no proj. & -- & 7.53e-2 & 0.72 [0.71, 0.73] & 0.00 \\
$u_c$ alone & -- & 3.87e-1 & 0.13 [0.13, 0.14] & 0.00 \\
\bottomrule\end{tabular}
\end{adjustbox}
\end{table}

The same construction applies to vector problems on the degree-of-freedom graph: the prolongations are $P_k\otimes I_2$ with the
Jacobi steps of \eqref{eq:jacobi} on the vector operator, which reproduce the rigid-body modes where the hypotheses of
Proposition~\ref{prop:constants} hold (\appref{app:impl}). We consider plane strain (Poisson ratio $0.3$) on the perforated square,
clamped on the outer boundary, traction-free on the holes, with a downward body force and a Young's modulus that is $100$ times
larger in four meandering bands than in the matrix; $L=3$ (7{,}356 unknowns) and $k_0=2$. The scalar island detector does not
apply, and the bands reach the clamped boundary, so OA-VCB is used without deflation, with its Galerkin output and otherwise the
configuration of the scalar problems (no tuning for elasticity). Its classical reference is the same Chebyshev V$(4,4)$ cycle with
Galerkin output; the contained fixed-interval Chebyshev member is slightly more accurate than it and the contained Jacobi member
$1.82$ times less accurate (Table~\ref{tab:contained}). This experiment has one data set with a held-out test split, no independent
data-seed replications and one OA-VCB training run. OA-VCB is more accurate than its reference on every test sample,
$G = 2.32$ [2.20, 2.44], and two-grid PCG first reaches its energy error after a median of $11$ iterations (range $4$--$32$;
Table~\ref{tab:elasticity}). The two operator-informed baselines, trained with three initialisation seeds each, are about four times
less accurate than the reference ($G = 0.24$--$0.27$). One preliminary case does not establish broad applicability to mechanics.

\section{Computing platforms and reproducibility}\label{app:platform}
\paragraph{Platforms.} Training, the stored evaluation records and the GPU inference times: NVIDIA T4 GPUs (Kaggle). Cost table
(Table~\ref{tab:cost}): one Intel Xeon processor at 2.1\,GHz, 2 threads. Stronger-smoother measurement (Table~\ref{tab:avt},
Figure~\ref{fig:avt}): a different machine, one Intel Xeon processor at 2.80\,GHz, 2 threads, on which OA-VCB was re-timed; its
CPU re-evaluation reproduces the stored OA-VCB errors to $5.6\times10^{-6}$ relative on channels and, on barriers, to within the
CPU/GPU round-off envelope documented in \appref{app:harmonic} (one sample at $1.4\times10^{-3}$, median $3\times10^{-6}$).
Timing scopes: classical rows report the mean over samples of the per-sample minimum over repetitions; OA-VCB's operator data
and inference report the minimum over repetitions of the whole-test-set time divided by the number of samples (inference in
batches of eight). Software: Python 3.13, PyTorch 2.14 (CPU), NumPy 2.5, SciPy 1.18 with OpenBLAS 0.3.34, PyAMG 5.3, Matplotlib
3.11; sparse direct solves with SuperLU.
\paragraph{Versions and regeneration.} Every training and evaluation record stores the hash \texttt{f8653612eb84}, an MD5
prefix of the top-level model and training sources computed by \path{run_identity.py}; it identifies the model code of every
stored record, not the diagnostic scripts in \path{diagnostics_cost/} and \path{mathcheck/}, the Kaggle drivers in
\path{kaggle_runs/}, or the table and figure generators. The code archive, available on request (see Code and data availability), identifies all of these by a SHA-256 list of every
file (\path{RELEASE_SHA256.txt}). Every table is regenerated from the stored per-sample records by \path{manuscript/make_tables.py},
which also writes \path{manuscript/tables/numbers.txt} (every number quoted in the text) and \path{manuscript/tables/MANIFEST.txt}
(the 91 input files used, as paths relative to the archive root); Figures~\ref{fig:results}--\ref{fig:avt} and
Figure~\ref{fig:setup} are regenerated by \path{manuscript/figures/make_figures.py}. The command sequence is given in the
archive's \path{README.md}.
